\section{Choosing initialization and the observer}
\label{sec:bb-initialization}

The preceding recorders supply complete instruction tables. We now
specify how their boxes are positioned and how the fixed material label
becomes the initial code. We compare one-time loading, a loading branch
repeated in every cycle, and chart translations that place the initial
code at the label itself. The latter two choices can preserve periodicity
from time zero. Each construction also needs a bound between integer
rule times; the endpoint coding alone does not give a material event.

Unless stated otherwise, the force in this section is the residual
\eqref{eq:bb-residual}, with pressure zero.  Its smoothness, all-order
mixed derivative bounds, finite effective prescription and uniqueness
are as in Proposition~\ref{prop:bb-realization}.  Whole-space forces
have one compact spatial support; torus forces have zero spatial mean.
All statements use a fixed positive computable viscosity, with the
relative interpretation of that proposition for arbitrary fixed
positive viscosity.

\subsection{Torus loading and direct chart placement}

\begin{proposition}[A move-first material test]\label{thm:bb-torus33}
The move-first recorder has a unit-torus realization, periodic after
$t=1$, whose fixed label $(1/4,1/2,1/4)$ enters
$\{x_1\in(1/2,1)\pmod1\}$ exactly when the work machine halts.
For a fixed machine only the loading interval depends on the input.
\end{proposition}
\begin{proof}\label{proof:bb-torus33}
Use Lemma~\ref{lem:bb-movefirst}, let $\Sigma$ be its full cell alphabet,
and enumerate its $N$ controls by $i=1,\ldots,N$.
Set $s=1/[16(N+1)]$ and give the $i$th control a square
of side $s$ centered at
\[
 m_i=(o_i+i/[4(N+1)],1/2),\qquad
 o_i=\begin{cases}5/8,&C_q\text{ with }q\text{ halting},\\1/8,&\text{otherwise}.
 \end{cases}
\]
Use odd digits and $B=2|\Sigma|+2$.  The code at height $z_0=1/4$
is $(m_i+s(E(L)-1/2,E(R)-1/2),z_0)$.
It is rational initially.  Lemma~\ref{lem:bb-onecell} gives $J$
branches with separately separated rectangles, all side lengths at
most $s$ and endpoints on the grid of mesh
$[32(N+1)B^2]^{-1}$.  For branch $j$ choose height
$z_j=1/2+j/[4(J+1)]$.  Use the three-phase lift of
Lemma~\ref{lem:bb-lift-parking}, with exponential middle scale
$\lambda_j^\theta$ and center interpolation.  A padding parameter
\[
 \epsilon=\min\{1/128,[320(N+1)B^2]^{-1},[40(J+1)]^{-1}\}
\]
separates the enlarged sheets in the source, height, and target phases
and keeps them inside the unit cube.  Their thickness may then be
chosen positive by Lemma~\ref{lem:bb-curl}.  For $J=0$ the repeated
field is zero.

During $[0,1]$ a localized curl translation takes the fixed label
along the segment to the initial code.  Both segment and a small
padding lie in the chart.  Repeat the branch pulse thereafter.
At each subsequent integer time the particle is the next table code.
If both endpoint controls are nonhalting, the interpolated center and
half-width bounds give $x_1<3/8+s<1/2$ throughout the motion.
The loader also stays below $1/2$ in a nonhalting run.  A halting
checkpoint lies in the upper half, reached in finitely many positive
local steps; an initial halt is reached at the end of loading.
This proves the event for every real time.
\end{proof}

\begin{proposition}[A guarded-window material test]\label{thm:bb-window37}
The table of Lemma~\ref{lem:bb-window-compiler} has a unit-torus
realization, periodic after $t=1$, with fixed label $(1/4,1/4,1/4)$
and event $x_1\in(1/2,7/8)\pmod1$, occurring exactly when the
work machine halts.  At work step $n$ its next
checkpoint occurs after exactly $2(p_n-h_n)+1\le4n+5$ unit periods.
\end{proposition}
\begin{proof}\label{proof:bb-window37}
Choose $r_0=2$ in Lemma~\ref{lem:bb-window-compiler}: initially the
frontier is at cell two, all other history cells are empty, and all
marks are zero. Thus $p_n=n+2$, with records in cells $2,\ldots,n+1$.
Split the full three-cell windows $[-1,1]$ and apply
Lemma~\ref{lem:bb-window-rectangles} with $(a,b)=(1,2)$.
Let $\mathcal A$ be the full cell alphabet of this window table and use
odd digits in base $B=2|\mathcal A|+1$.
For $S$ controls let $\epsilon=1/[16(S+1)]$ and place their squares
at $o_s+\epsilon[0,1]^2$, where
\[
 o_s=(\xi_s,1/4+2\kappa(s)\epsilon),\quad
 \kappa(s)\in\{1,\ldots,S\},\quad
 \xi_s=\begin{cases}5/8,&s\text{ halting Ready},\\1/4,&\text{otherwise}.
 \end{cases}
\]
Write $J$ for the number of branches and $e_j$ for branch $j$'s cursor
displacement. The full source and target rectangle gaps are at least
$\epsilon B^{-3}$.  Put their common height at $1/4$, assign branch
$j$ height $1/2+j/[4(J+1)]$, and use the exponential middle scale
$\operatorname{diag}(B^{-e_j\theta},B^{e_j\theta},1)$.
Take one tenth of
$\min(1/32,\epsilon B^{-3},1/[4(J+1)])$ as padding.
Horizontal source and target gaps and the private heights separate
the three phases, while all widths remain at most $\epsilon$.
The strain can equivalently be generated by
$(0,0,-e_j\dot\theta\log B\,(x_1-c_{j1})(x_2-c_{j2}))$,
retaining the explicit diagonal-strain implementation.

A first unit interval loads the label along the segment to the rational
initial code.  Thereafter the exact rectangle action and the compiler
lemma give checkpoint times $t_0=1$ and
$t_{n+1}=t_n+2(p_n-h_n)+1$.  A nonhalting center stays in
$[1/4,1/4+\epsilon]$, with half-width at most $\epsilon/2$;
thus every intermediate first coordinate remains below $1/2$.
The loader has the same exclusion.  Halt codes are in
$[5/8,5/8+\epsilon]$, inside the observer, including an initial halt.
The history block starts at the known absolute cell two, so the
head-relative code also recovers absolute indexing at checkpoints.
\end{proof}

\begin{proposition}[Two direction-recording placements]
\label{thm:bb-aligned38}\label{thm:bb-direction39}
The direction recorder gives both of the following realizations.
\begin{enumerate}
\item On $\R^3$, with its passive origin bit, there is a force periodic
from time zero and a fixed-origin event $X_1(t;0)<-5$.  Input dependence
can be confined to two translations of the state charts.  On every
implemented full box the first coordinate is a convex combination
of its endpoint first coordinates.
\item On the unit torus, without the passive track, a force periodic
after $t=1$ has fixed label $(1/8,1/8,1/4)$ and event
$x_1\in(1/2,7/8)\pmod1$.  For a fixed machine its repeated part is
independent of the input.
\end{enumerate}
Both events are equivalent to halting.
\end{proposition}
\begin{proof}\label{proof:bb-aligned38}\label{proof:bb-direction39}
For the first case introduce a fresh terminal state $\dagger$ and
replace every missing instruction and every originally terminal
state--symbol pair by a write-preserving stationary transition into it.
The original initial control is therefore nonterminal in this normalized
table, even when the original machine is initially halted. Apply
Lemma~\ref{lem:bb-direction} with its passive origin track, and let
$\mathcal A$ be the resulting full cell alphabet. Use odd digits in
base $B=2|\mathcal A|+1$
and initial fractions $x_*,y_*$.  Index the initial main state by zero,
the main $\dagger$ states by distinct negative integers, and the rest
by distinct positive integers.  State $s$ has placement
\[
 (x,y)\longmapsto(10k(s)+x-x_*,y-y_*).
\]
The initialized point is zero; nonhalt codes have first coordinate
$>-1$, and halt codes $<-5$.  Thicken the rectangles of
Lemma~\ref{lem:bb-onecell} by $[-1,1]$, give branch $i$ height $10i$,
and use the linear scale $g_i=1+\theta(\lambda_i-1)$ in the lift
schedule.  A padding smaller than one tenth of the minimum of one
and all within-family planar gaps suffices.  The middle strain has
potential $(0,0,(\dot g_i/g_i)(x_1-c_{i1})(x_2-c_{i2}))$.
Equation~\eqref{eq:bb-convex-coordinate} excludes every nonhalting
visit to $x_1<-5$.  The table and routing are fixed for the machine;
only $x_*,y_*$ change with the input.  The passive origin bit remains
at physical cell zero, so no absolute tape information is discarded.

For the second case first merge all halting work states into a single
state $q_h$, redirecting incoming transitions and retaining initial halting.
Delete the passive track from the direction recorder, and now let
$\mathcal A$ denote that full alphabet. Use odd digits with
$B=2|\mathcal A|+2$ and
$\lambda=1/[16(J+1)]$, where $J$ counts the controls.  Number them
from zero and set
\[
 S_s(x,y)=(A_s+\lambda x,1/4+\lambda y),\qquad
 A_s=2\lambda\operatorname{index}(s)+
 \begin{cases}5/8,&s=M(q_h,e),\\1/4,&\text{otherwise}.
 \end{cases}
\]
The nonhalt squares lie in $[1/4,3/8]\times[1/4,1/4+\lambda]$;
halt squares lie in $[5/8,3/4]\times[1/4,1/4+\lambda]$.
Write $m$ for the number of branches. Use coding height $1/4$,
branch heights $1/2+i/[4(m+1)]$, and
the same linear first scale.  One hundredth of the minimum of
$1/16$, $1/[4(m+1)]$, and the source and target planar gaps is
valid padding; absent pair gaps are omitted.  All supports lie inside
the cube, and all horizontal widths are at most $\lambda$.
A localized translation with padding $1/1000$ loads the rational
initial point from the fixed label in one unit interval.

In a nonhalting period the center has first coordinate at most $3/8$
and half-width at most $\lambda/2$, hence is below $1/2$; loading
has the same bound.  Halting codes lie in the observer.  The exact
sampling induction gives one table step per period, with checkpoint
cycle $2(k+1-h_k)+3$.  The record block beginning at cell one
recovers absolute indexing at checkpoints.  An initial halt is
observed at time one.
\end{proof}


\subsection{A one-time torus loader}

The next placement uses a separate first interval to load the code. Only
the instruction cycle is then repeated, so periodicity begins at time one.

\begin{theorem}[A separate one-unit loader on the torus]
\label{ba:loaded-torus}
For every machine and finite input and every fixed computable $\nu>0$,
there is an effective smooth
mean-zero force, one-periodic for $t\ge1$ with bounded mixed derivatives,
whose unique smooth zero-data solution with zero normalized pressure
satisfies
\[
 \exists t\ge0:\ \Phi_t(1/2,1/2,1/4)
       \in(1/16,1/4)\times\mathbb T^2
 \quad\Longleftrightarrow\quad\text{halting}.
\]
Uniqueness holds in the classical class of Section~\ref{sec:model}.
\end{theorem}
\begin{proof}\label{ba:loaded-torus-proof}
Use the single-halt move-first recorder of
Lemma~\ref{lem:bb-movefirst}, without an entry state.  Let
$\Gamma$ be its full alphabet and $N$ its number of controls.  Put
$B=2|\Gamma|+1$, $a=1/(16N)$, the halt square at
$(1/8,1/4)+a[0,1]^2$, and the other squares at
$(1/2+2ja,1/4)+a[0,1]^2$, $0\le j\le N-2$.
Use the minimal subdivision, so right and stay moves are not split by
the left letter.  Its gaps are at least $a/B^2$.
For $M$ branches choose
\[
 h=\frac1{64(M+1)},\quad Z_i=\frac12+\frac{i}{4(M+1)},\quad
 \varepsilon=\min\{1/128,h/4,a/(4B^2)\}.
\]
Apply linear-first-scale full-box routing from height $z_0=1/4$.
Middle heights differ by $16h$; source and target gaps handle the
vertical phases.  Centers lie in
$[1/8,5/8]\times[1/4,5/16]$, widths are at most $a$, and all padded
boxes lie in $(1/16,15/16)^3$.  This is a construction on the specified
unit torus, with no change of physical viscosity.

In the first unit interval translate a box of side $1/64$ from the
fixed label to the rational initial code, using Lemma~\ref{lem:bb-curl}.
The segment and a small padding lie inside the same chart.  Repeat
only the instruction cycle afterward.  Code induction proves the
checkpoint identities.  A halting code has first coordinate in
$[1/8,3/16]$, including an initial halt reached by loading.  For a
nonhalting run both centers of every used branch have first coordinate
at least $1/2$, with half-width at most $a/2$; loading has the same
lower center bound.  Hence the particle always stays to the right of
$1/4$.  Residual realization proves every analytic assertion.  The
separate loading field accounts for the qualification $t\ge1$ in the
periodicity statement.
\end{proof}


\subsection{Repeated torus loading branches}\label{ba:torus}

A loader can instead be one branch of every period. To keep its target
separate from all computational targets, place the initial code in a state
square with no incoming instruction. This condition concerns the entire
target family, not only the distinguished execution. We implement the
same startup principle first with full boxes, then with fixed columns.

\begin{theorem}[A loader repeated as part of the processor]
\label{ba:periodic-torus}
For every finite machine and input, and every fixed computable $\nu>0$,
there is an effective smooth mean-zero force on the unit flat torus,
one-periodic from time zero, with all mixed derivatives bounded.  Its
zero-data solution has pressure zero, is unique in the classical class
of Section~\ref{sec:model}, and satisfies
\[
 \exists t\ge0:\ \Phi_t(1/8,1/8,1/4)
             \in\{5/8<X_1<7/8\}
 \quad\Longleftrightarrow\quad\text{the machine halts}.
\]
The one repeated cycle implements all instruction branches and its
loading branch.
\end{theorem}
\begin{proof}\label{ba:periodic-torus-proof}
Normalize to one halt state and prepend a new nonhalting state $q_I$
whose rule preserves the scanned letter, stays, and enters the original
initial state.  No machine rule targets $q_I$.  Apply the move-first
recorder of Lemma~\ref{lem:bb-movefirst} and omit the unused controls
$P_{q_I},T_{q_I}$. Then the entry control $C_{q_I}$ has no incoming local
instruction.  Write $\Gamma$
for this recorder's full alphabet.  Put $B=2|\Gamma|+1$, let $N$
be the number of controls, and set
$a=1/(16(N+1))$.  Use the halt square
$(3/4,1/4)+a[0,1]^2$ and nonhalting squares
$(1/4+2al,1/4)+a[0,1]^2$, $0\le l\le N-2$.
The full subdivision of Lemma~\ref{ba:rectangles} gives source
half-widths $a/(2B)$ in both coordinates.

Add a square with those half-widths centered at $(1/8,1/8)$ and a
translation to the square centered at the initial entry code.  Each
initial tail is at distance at least $1/(B-1)>1/(2B)$ from the endpoints
of $[0,1]$, so the new target lies strictly inside the entry square.
There is no other target in that square.  Both enlarged branch lists
therefore remain separately separated, with gaps at least $a/B^2$.
The extra source is separated from all state squares.

Number the branches, including loading, by $1\le i\le M$.  Set
\[
 z_0=\tfrac14,\qquad h=\frac1{16(M+1)},\qquad
 Z_i=\tfrac12+\frac{i}{4(M+1)},\qquad
 \varepsilon=\frac1{10}\min(h,a/B^2).
\]
Thicken each planar source by $[z_0-h,z_0+h]$.  Lift the resulting
full box to center height $Z_i$, move it horizontally, and lower it
to $z_0$.  For an instruction with displacement $d_i$, use
$D_i=\operatorname{diag}(B^{-d_iw_2},B^{d_iw_2},1)$;
loading has $d_i=0$.  Its center is the planar endpoint interpolation
plus $(Z_i-z_0)(w_1-w_3)e_3$.  Adjacent middle heights are $4h$ apart,
while the padded half-thickness is $h+\varepsilon<2h$.
Source and target gaps handle the other phases.  Horizontal half-widths
never exceed $a/2$, horizontal centers lie between $1/8$ and $3/4+a$,
and vertical centers lie between $1/4$ and $3/4$.  The padded motions
are thus strictly inside the unit cube.  Lemma~\ref{lem:bb-curl} gives a
chart field, extended periodically in space and time.  Its zero time
collars give zero initial velocity.

The first cycle loads the fixed particle.  Each subsequent cycle maps
its exact code to the next local code until the first halt, or forever
in a nonhalting run.  Lemma~\ref{lem:bb-movefirst} identifies the finite
checkpoint times with machine steps.  The added entry state handles an
original initial halt.  A halt code is in the target slab.  For a
nonhalting execution every used source and target center, including the
loader, has first coordinate at most $3/8$.  Exponential interpolation
keeps the center below that bound and the half-width below $a/2$;
the particle cannot reach $5/8$ at any intermediate time.  The final
transition into a halt is not subject to this nonhalting-endpoint bound.
Lemma~\ref{ba:realization} gives the force, its mean, its effective
derivatives, and uniqueness.  Repeating the loading field is harmless
to the initialized trajectory and makes the force periodic from zero.
\end{proof}

\paragraph{A fixed-column implementation.}\label{bc:unit-periodic-section}
The same untargeted-entry condition also permits a repeated initializer
using fixed source and target columns. The next construction retains its
explicit column potentials and its fixed label; its middle motion uses
linear interpolation of the first planar scale.

\begin{theorem}[A periodic initializer on the unit torus]
\label{bc:unit-periodic-event}
There is an effective halting realization on $\T^3$ with fixed label
$a_*=(3/8,1/4,1/2)$ and fixed strip $5/8<x_1<7/8$, using a smooth
mean-zero general force of period one from time zero. Its initial velocity
and prescribed pressure are zero, every mixed derivative is bounded, and
its kinetic energy is uniformly bounded. The solution is unique in the
torus comparison class.
\end{theorem}
\begin{proof}\label{bc:unit-periodic-event-proof}
Normalize to a single halt and add a new nonhalting start $q_s$, with no
incoming original transition, which performs a dummy unchanged stay move
to the old initial state. This preserves even initial halting. Use the
move-first recorder of Lemma~\ref{lem:bb-movefirst}, with the control names
of Section~\ref{bc:section}, and retain its frontier-writing and return
controls $C_q,D_q$
only for $q$ in the set $Q^+$ of targets of the normalized work table,
including the new dummy transition. In particular its destination, the
old initial state, belongs to $Q^+$, whereas $q_s\notin Q^+$. No recorder
rule therefore targets the ready control $W_{q_s}$. Its
initial frontier is still two and the simulation proof is unchanged.

Use odd digits in base $B=2K+2$, where $K$ is the recorder alphabet size.
If there are $n$ nonterminal controls, put $\tau=1/[16(n+1)]$. Give
nonterminal states squares of side $\tau$ at lower corners
$(1/8+2\tau j,1/4)$, $0\le j<n$, and put the terminal square at
$(3/4,1/4)$. All nonterminal first coordinates are below $1/4$.
Refine every rule, including right and stationary rules, by the left
letter $\ell$. The full source is $I_\ell\times I_\gamma$, and the
three image rectangles are respectively
\begin{equation}\label{bc:unit-rectangles}
 I_{\beta\ell}\times[0,1],\qquad
 I_\ell\times I_\beta,\qquad [0,1]\times I_{\ell\beta}.
\end{equation}
They follow from the same affine formulas as Lemma~\ref{lem:bb-onecell}.
The refined sources are separately disjoint; the incoming displacement
and written symbol, together with the extra left prefix where applicable,
separate the images. Thus full geometric separation still holds.

Let $w_*=(3/8,1/4)$ and let $w_{\rm in}$ be the rational input code in
the untargeted start square. Add the translation from
$w_*+[-\delta,\delta]^2$ to $w_{\rm in}+[-\delta,\delta]^2$, where
$\delta=\tau/(10B^2)$. The source lies away from every state square.
The target lies strictly inside the start square: odd-digit codes have
distance at least $1/(B-1)$ from zero and at least $2/(B-1)$ from one,
both larger than $1/(10B^2)$. It meets no computational image because
none targets that state. This extended list still has separately disjoint
source and target families, even though a target can overlap a source.

We give the torus transport explicitly to keep the startup claim separate
from the whole-space column construction. All rectangles initially lie in
$z_0=1/2$. For $m$ branches assign heights
$z_j=1/2+j/[4(m+1)]$. Choose planar source and target cutoffs equal to
one on neighborhoods of their rectangles, with padding less than one tenth
of every relevant pairwise gap and chart-face margin; a missing pairwise
minimum is omitted. Let $\eta(z)$ equal one near $[1/2,3/4]$ and have
padding $1/16$. The source column field
\[
 Z_j=\nabla\times(0,x\chi_j(x,y)\eta(z),0)
\]
equals $e_3$ along its source column; the target analogue $Z'_j$ does
so along the target column. These fields have compact support in the
open unit cube and zero spatial mean after periodization.

For the middle phase use the center path
$c_j=(1-\mu)p_j+\mu q_j$ and scale
$l_j=1-\mu+\mu a_j$, where $a_j$ is the branch's horizontal scale.
Choose $\xi_j$ equal to one near the planar coordinate bounding rectangle
of the endpoints, and disjoint height cutoffs $\rho_j$ at $z_j$ with
padding $1/[16(m+1)]$. For three successive flat ramps $\theta_1,
\theta_2,\theta_3$, put $c_j=c_j(\theta_2)$, $l_j=l_j(\theta_2)$ and
\[
 \Psi_j=\xi_j\rho_j\left[
 \dot c_{j,1}(y-c_{j,2})-\dot c_{j,2}(x-c_{j,1})
       +\frac{\dot l_j}{l_j}(x-c_{j,1})(y-c_{j,2})\right].
\]
The full velocity is
\begin{equation}\label{bc:unit-V}
 V=\sum_j(z_j-z_0)\theta_1'Z_j
       +\sum_j(\partial_y\Psi_j,-\partial_x\Psi_j,0)
       +\sum_j(z_0-z_j)\theta_3'Z'_j.
\end{equation}
It lifts all full rectangles, performs the reciprocal affine motions in
separate layers, and lowers them. In the middle phase the exact planar path
is $c_j+\diag(l_j,l_j^{-1})(\zeta-p_j)$. The swept-rectangle bound in
Lemma~\ref{ba:two-interpolations} keeps it inside the coordinate bounding
rectangle of the endpoints. The first coordinate is an endpoint convex combination.
The cutoffs are therefore one along every claimed path. Differentiation
and uniqueness prove the assigned full maps. The entire construction is
supported in the open unit cube and has zero collars in phase.

Extend it with period one in phase and periodically in space. It is zero
at time zero, so its residual has the required period from time zero.
During the first period the distinguished particle follows the initializer;
its first coordinate stays below $1/2$. Thereafter it follows the recorder.
Every nonterminal transition remains below $1/4$, whereas a terminal code
lies inside $5/8<x_1<7/8$. It cannot return to the initializer source,
whose square lies at $x_1=3/8$. The dummy instruction handles initial
halting without changing the observation rule. Lemma~\ref{bc:analytic}
now proves all analytic assertions, including uniform energy on the compact
torus. The force formula includes every branch, the initializer included,
in each period; it uses no precomputed execution.
\end{proof}


\subsection{A repeated loading instruction}

The next construction incorporates loading in a different symbolic table:
its instructions replace finite prefixes of two stacks. A replacement can
change planar area, so the third coordinate compensates for the product
of the two planar scales. Evacuation and delivery then realize the
resulting volume-preserving boxes in every period.

\begin{theorem}[Periodic prefix-box realization]\label{thm:bb-prefix29}
For every deterministic machine and finite input, a compactly supported
smooth force on $\R^3$, periodic with period one from time zero,
realizes the fixed-origin halting event
\[
 \exists t\ge0:\quad X(t;0)\in(-4,-1)\times(-1,2)\times(-1,1)
 \quad\Longleftrightarrow\quad\text{halting}.
\]
The force has bounded mixed derivatives and the unique zero-data
solution in the class of Proposition~\ref{prop:bb-realization}.
Its finite program uses full volume-preserving boxes; no solenoidal
condition is imposed on the force itself.
\end{theorem}
\begin{proof}\label{proof:bb-prefix29}
Normalize to one halting work state $q_h$, redirecting incoming
transitions to it and replacing missing nonhalting rules by unchanged-letter
stay moves into it. This preserves initial halting.
Write $\Gamma$ for the work alphabet, $b$ for its blank, $q_0$ for the
initial control, and $\omega$ for the finite input word.
For top-first stacks $L,R$, a rule
$(s,\alpha,\beta)\rightsquigarrow(r,\gamma,\eta)$ replaces the
two displayed prefixes and leaves their independent infinite tails
unchanged.  Introduce fresh state $s_0$, boundary letter $S$ and mark
letter $M$.  The preliminary rules are
\[
 (s_0,\varepsilon,\varepsilon)\rightsquigarrow(q_0,S,\omega)
\]
and, for each work transition $(q,a)\mapsto(r,c,d)$,
\[
\begin{array}{c|l}
 d&\text{preliminary rules}\\\hline
1&(q,\varepsilon,a)\rightsquigarrow(r,c,\varepsilon)\\
0&(q,\varepsilon,a)\rightsquigarrow(r,\varepsilon,c)\\
-1&(q,v,a)\rightsquigarrow(r,\varepsilon,vc)\quad(v\in\Gamma),\\
  &(q,S,a)\rightsquigarrow(r,S,bc).
\end{array}
\]
The final row exposes the implicit blank beyond the left boundary.
Index each preliminary occurrence by $j$, with target $r(j)$,
and introduce a private control $g_j$ and record letter $\ell_j$.
Replace that occurrence by
$(s,\alpha,\beta)\to(g_j,\gamma,M\eta)$ and add
\begin{equation}\label{eq:bb-prefix-table}
\begin{aligned}
 (g_j,v,\varepsilon)&\to(g_j,\varepsilon,v)&&v\in\Gamma,\\
 (g_j,S,\varepsilon)&\to(k_{r(j)},S\ell_j,\varepsilon),\\
 (k_q,\varepsilon,v)&\to(k_q,v,\varepsilon)&&v\in\Gamma,\\
 (k_q,\varepsilon,M)&\to(q,\varepsilon,\varepsilon).
\end{aligned}
\end{equation}
Sources are separated by control and tested top symbols.  Into $g_j$,
entry has right top $M$, whereas a loop has a work letter there.
Into $k_q$, entry has left prefix $S\ell_j$, whereas a loop has a
work letter at left top; distinct entries have distinct $\ell_j$.
Into a work state only the last rule enters.  Thus both source and
target cylinders are disjoint for arbitrary tails.

After a preliminary replacement the stacks have form $wSB,MY$,
with $w\in\Gamma^*$.  The forward loop transfers $w$, leaving
$SB,w^{\rm rev}MY$.  The boundary rule inserts $\ell_j$ behind
$S$, the reverse loop restores $w$, and the final rule removes $M$.
After $2|w|+2$ further instructions the stacks are
$wS\ell_jB,Y$ at control $r(j)$.  At a work control, $w$ lists
represented cells left of the head, nearest first; $R$ lists the
current cell and right tail, and cells beyond $S$ on the left are
implicit blanks.  The preliminary rules perform exactly the work
write and move.  Starting at $(s_0,b^\infty,b^\infty)$, the first
work visit is at rule count three, and consecutive work visits differ
by $3+2|w_{\rm new}|$.  This proves finite continuation and the
halting equivalence, including an initially halted input.  The
retained records also recover all head displacements.

Enumerate the enlarged alphabet, blank first, by even digits
$0,2,\ldots,2m-2$ in base $K=2m$.  Use prefix intervals $I(v)$
and codes $E$ as in Lemma~\ref{lem:bb-onecell}.  Set offsets
$o_{q_h}=-3$, $o_{s_0}=0$, and $3,6,9,\ldots$ for the other controls.
Encode $(s,L,R)$ by $(o_s+E(L),E(R),0)$; its initial value is zero.
For a rule $(s,\alpha,\beta)\to(s',\gamma,\eta)$ put
\[
 a_i=K^{|\alpha|-|\gamma|},\quad
 b_i=K^{|\beta|-|\eta|},\quad c_i=(a_ib_i)^{-1}.
\]
The full source box is
$(o_s+I(\alpha))\times I(\beta)\times[-1,1]$;
the target is
$(o_{s'}+I(\gamma))\times I(\eta)\times[-c_i,c_i]$.
The center-to-center map has matrix $\operatorname{diag}(a_i,b_i,c_i)$
and determinant one.  It implements the prefixes on all codes.
Disjoint cylinders give separated full planar projections in each
family, including blank-tail endpoints.

Park box $i$ at $(4i,4,0)$, reserving the horizontal rectangle
$[4i-1/2,4i+1/2]\times[7/2,9/2]$.
Let $R=\max(1,c_1,\ldots,c_N)$ and travel at height $Z=2R+3$.
Evacuate all sources to their parking sites, reshape there by
$\operatorname{diag}(a_i^s,b_i^s,c_i^s)$, then deliver all targets.
The two horizontal widths remain at most one, and vertical half-widths
at most $R$.  Let $g$ be the minimum of one and the within-family
gaps in the source-plus-parking and target-plus-parking projections.
Padding smaller than $g/8$ works during vertical motion and strain;
at travel height the moving bottom is at least $Z-R>R+2$.
Lemma~\ref{lem:bb-lift-parking} gives the full $7N$-motion schedule.

Repeat this unit pulse from zero.  Nonhalting endpoint boxes and all
parking boxes have nonnegative first coordinate.  Translations preserve
that bound, and strain takes place in positive parking envelopes.
Therefore a nonhalting path never enters the negative observation box.
A halting code lies in $[-3,-2]\times[0,1]\times\{0\}$ inside it.
The residual proposition finishes the proof.  Loading was a repeated
instruction, so no exceptional first-time force is required.
\end{proof}


\begin{remark}[A different observation regime]
For comparison, the distinct velocity-field detector of
\cite[Theorem~1.1(2) and Section~5]{OpenAI379EulerianC} uses
$\R^2\times\T$ and the event
$\int_{\{X_2>0\}\times\T}u_3>1/2$. Its force has globally bounded
mixed derivatives and horizontal support compact on each finite time
interval. The nonnegative scalar velocity has diffusion tails; after
loading its total mass is one. Its cutoff estimate gives loss
$\delta<1/24$, nonhalting mass at most $\delta$ in the observed half-space,
and halting mass above $3/4$. Those conclusions follow from the cited
advection--diffusion proof and have no uniform spatial support assertion.
They are a separate force and observation regime from the material test
proved here.
\end{remark}
